% Historical calculation retained for provenance; not included by either active edition.
\subsection{Exact direct output formulas}
\label{subsec:new-direct-output-formulas}

The following outgoing capacities follow from
Proposition~\ref{prop:new-exact-local-set} and supply the CE1 return.

The next formulas quantify the boundary--interior tradeoff along a path.
A large radial demand can make a large incoming boundary reach incompatible
with a large outgoing one. Only the selected admissible branch may be used:
a formal root of a quadratic is not automatically the boundary capacity.
The threshold and return arguments that follow use precisely these outputs.

For a nonsupercritical triangle with incoming reach at least $a$ and radial
reach at least $c$, let
\[
 \mathcal B_c(a)=
 \max\{b:(a,b,c)\text{ is admissible and }a+b\le1\}.
\]
This is the outgoing capacity of a single triangle.

For $1/2<c\le\sqrt3/2$, put
\[
 h_c=\frac{2c}{1+\sqrt{16c^2-3}}.
\]
For $\sqrt3/2<c<1$, put
\[
 \ell(c)=\frac c2\left(1-\sqrt{4c^2-3}\right),
 \qquad r(c)=c-\ell(c).
\]
Also set
\[
 \mathcal Q_-(a,c)=\frac{\sqrt{a^2-ac+c^2}}c-a,
\]
\[
 \mathcal Q_+(a,c)=
 \frac{2(1-a^2)c^2}
 {2ac^2+c+
 \sqrt{(2ac^2+c)^2-4(1-c^2)(1-a^2)c^2}}.
\]

\begin{proposition}[Four direct high-radial outputs]
\label{prop:new-four-direct-outputs}
For $1/2<c\le\sqrt3/2$,
\[
 \mathcal B_c(a)=
 \begin{cases}
  1-a,&0\le a\le1-c,\\
  \mathcal Q_+(a,c),&1-c<a<h_c,\\
  h_c,&a=h_c,\\
  \mathcal Q_-(a,c),&h_c<a<c,\\
  1-a,&c\le a\le1.
 \end{cases}
\]
For $\sqrt3/2<c<1$,
\[
 \mathcal B_c(a)=
 \begin{cases}
  1-a,&0\le a\le1-c,\\
  \mathcal Q_+(a,c),&1-c<a\le\ell(c),\\
  \ell(c),&\ell(c)<a<r(c),\\
  \mathcal Q_-(a,c),&r(c)\le a<c,\\
  1-a,&c\le a\le1.
 \end{cases}
\]
At $a=\ell(c)$, the selected value is $r(c)$; immediately to the right it is
$\ell(c)$.
\end{proposition}

\begin{proof}
Since $c>1/2$, any containing V triangle contains the radial midpoint.
Therefore $a+b\le1$, and the supercritical cell of Proposition
\ref{prop:new-exact-local-set} is absent.  On the cap $b=1-a$, the $T$-cell
inequality becomes
\[
 M(c-M)\le0,
\]
where $M=\max\{a,b\}$.  Thus the linear cap is feasible exactly when
$a\le1-c$ or $a\ge c$.

Off the linear cap, a maximal point lies on $F_L=0$ or $F_T=0$.  The selected
$L$-cell equation has the two ordered roots $\ell(c),r(c)$; the selector
$q\le0$ retains the smaller coordinate.  The ordered halves of $F_T=0$ give
$\mathcal Q_+$ and $\mathcal Q_-$.  Their equality boundary $a=b$ is $h_c$ in the low-radial
range, while their $q=0$ boundaries are
$(\ell(c),r(c))$ and $(r(c),\ell(c))$ in the high range.  These conditions
produce the two displayed formulas.

The formal other root in the $\mathcal Q_+$ quadratic violates the geometric selector
$c\le2\max\{a,b\}$ and is discarded.  At the high-range junction, the
selected $T$-cell point has value $r(c)$, whereas the adjacent selected
$L$-cell point has value $\ell(c)$; this is the genuine downward jump stated
above.
\end{proof}

\subsection{High-radial and selected-chord calculations}

For
\[
 0<e<1-\frac{\sqrt3}{2},
\]
define the selected low root
\begin{equation}
 \epsilon(e)=\frac{1-e}{2}
 \left(1-\sqrt{4(1-e)^2-3}\right).
\label{eq:selected-low-root}
\end{equation}
It satisfies
\begin{equation}
 e<\epsilon(e),\qquad
 \epsilon(e)<\frac{2e}{1-2e},
\label{eq:selected-low-root-basic-bounds}
\end{equation}
and, for $e\le1/8$,
\begin{equation}
 \epsilon(e)\le2e+5e^2.
\label{eq:selected-low-root-small-bound}
\end{equation}
These follow by substituting the proposed upper bounds in the selected
quadratic and checking its sign.

\begin{lemma}[Direct high-radial threshold]
\label{lem:new-direct-threshold}
Let $0<e<1-\sqrt3/2$. A nonsupercritical V triangle has incoming boundary reach $A$ and
outgoing boundary reach $B$.  If
\[
 C\ge1-e,
 \qquad
 A>\epsilon(e),
\]
then
\[
 \boxed{B\le\epsilon(e).}
\]
\end{lemma}

\begin{proof}
The condition $C>1/2$ forces $A+B\le1$.  In the exact cells of
Proposition~\ref{prop:new-exact-local-set}, the only selected component with
$A>\epsilon(e)$ has $B\le\epsilon(e)$; equivalently the two roots of the
selected $T$-cell quadratic separate the feasible fiber.  The selector
$c\le2\max\{A,B\}$ removes the formal high root.
\end{proof}

\begin{lemma}[Selected-arc chord bounds]
\label{lem:new-selected-chords}
Put $c=1-e$, with $0<e<1/4$, and define
$q_*(p)=1-\mathcal Q_+(p,c)$ on the selected branch. Its right endpoint is
\[
 (p_e,q_e)=
 \begin{cases}
 (\epsilon(e),e+\epsilon(e)),&c>h,\\
 (h_c,1-h_c),&c\le h,
 \end{cases}
 \qquad h=\sqrt3/2.
\]
The function $q_*$ is increasing and concave on $[e,p_e]$, with
$q_*(e)=e$. Thus, for $e<p\le p_e$,
\[
 q_*(p)\ge p+\lambda_e(p-e),\qquad
 \lambda_e=\frac{q_e-p_e}{p_e-e}.
\]
One has $\lambda_e>1-4e$ when $0<e<1-h$, and
$\lambda_e>1-5e$ throughout $0<e<1/4$.
A following actual incoming reach $q$, forced by a weak boundary handoff,
satisfies $q\ge q_*(p)$; equality is not assumed for actual reaches.
\end{lemma}
\begin{proof}
The selected boundary is parametrized by
\[
 q_*=e+cx,\qquad p=e+cx-1+\sqrt{1-x+x^2}.
\]
Indeed its selected quadratic gives
$x(1-x)=(q_*-p)(2-q_*+p)$. If $p=p(x)$, then
\[
 p'(x)=c+\frac{2x-1}{2\sqrt{1-x+x^2}}\ge c-\frac12>0,
 \qquad p''(x)=\frac3{4(1-x+x^2)^{3/2}}>0
\]
on the selected interval $0\le x\le1$.
Hence $dq_*/dp>0$ and $d^2q_*/dp^2=-cp''/(p')^3<0$.
The endpoint formulas follow from \zcref{prop:new-four-direct-outputs},
and concavity places the graph above its endpoint chord.
For $c>h$,
\[
 \lambda_e=\frac e{\epsilon(e)-e}
 >\frac{1-2e}{1+2e}>1-4e,
\]
using \eqref{eq:selected-low-root-basic-bounds}.
For $3/4<c\le h$, the formula for $h_c$ gives
\[
 h_c\le\frac3{2(1+\sqrt6)}<\frac7{16},\qquad
 \lambda_e>\frac2{7-16e}>1-5e.
\]
For the first bound, $h_c$ decreases with $c$ on this interval and
$17^2<49\cdot6$. For the second, $e\ge1-h>1/8$, and
$80e^2-51e+5<0$ on $[1/8,1/4]$: it decreases there and has value
$-1/8$ at $1/8$.
Finally the preceding outgoing reach is at most $\mathcal Q_+(p,c)$;
a weak handoff therefore gives $q\ge1-\mathcal Q_+(p,c)$.
\end{proof}

\subsection{CE1 reverse-path calculation}

\label{subsec:new-ce1-direct-certificate}

The first lemma uses only the $T_4$ radial demand. The full return adds
the $T_3$ and $T_2$ demands; BC establishes them before using the latter.

Use the signed center parameters of \zcref{prop:signed-center-normal-form},
with $w=1-R$, $A=\alpha$, and $D=\delta$.
Assume
\begin{equation}
 A>0,
 \qquad D>0,
 \qquad A+wD<P,
 \qquad D+RA\ge P.
 \label{eq:ce1-center-assumptions}
\end{equation}
Put
\begin{equation}
 X=R-D,
 \qquad m=\frac AR,
 \qquad h_0=\frac{\eta+A+D}{2R}.
 \label{eq:ce1-auxiliary-parameters}
\end{equation}

\begin{lemma}[First CE1 return step]
\label{lem:ce1-first-return-step}
Assume \eqref{eq:ce1-center-assumptions}, and let
$T_1,\ldots,T_4$ be nonsupercritical V triangles of arbitrary types,
with actual reaches satisfying
\begin{equation}
 A_1\ge X,\qquad B_i+A_{i+1}\ge1\ (1\le i\le3),\qquad B_4\ge h_0.
 \label{eq:ce1-b4-floor}
\end{equation}
If $C_4\ge1-A$, then the reflected $T_4$ state lies on the selected
$\mathcal Q_+$ component and
\begin{equation}
 B_3>L_1,\qquad L_1=(2-4A)h_0-(1-4A)A.
 \label{eq:ce1-first-chord}
\end{equation}
No own-radial demand at $T_2$ or $T_3$ is assumed.
\end{lemma}

\begin{proof}
The sign difference in \zcref{eq:ce1-center-assumptions} gives
\[
 RD>wA,
\]
so the center exit on $r_3$ is $m=A/R$.  Moreover
\begin{equation}
 0<A<P<\frac18,
 \qquad 0<D<\frac R2,
 \qquad 0<m<\frac w2<\frac12.
 \label{eq:ce1-small-parameters}
\end{equation}
We first record
\begin{equation}
 \epsilon(A)<X,
 \qquad h_0>A.
 \label{eq:ce1-root-height}
\end{equation}
Indeed, $A+wD<P$ and $D=R-X$ give $A<wX-\eta$.  Convexity of
\[
 wX-\frac{X}{2(1+X)}
\]
on $R/2<X<R$ gives
\[
 A<\frac{X}{2(1+X)}.
\]
The selected-root bound
\zcref{eq:selected-low-root-basic-bounds} yields $\epsilon(A)<X$.  Also
\[
 2R(h_0-A)=\eta+D+(1-2R)A>0;
\]
when $R>1/2$, use $A<P=\eta E$ for the last sign.

The assumed transverse radial reach at $T_4$ is at least $1-A$.  Apply the
exact local catalogue to the reflected boundary pair $(B_4,A_4)$.  If the
local point is not on the selected $\mathcal Q_+$ component, the constant, $\mathcal Q_-$,
and linear components together with \zcref{eq:ce1-root-height} give
\[
 1-A_4\ge1-\epsilon(A)>1-X.
\]
The center-free edges then carry this lower bound backward to
$B_1>1-X$, contradicting $A_1\ge X$ and $A_1+B_1\le1$.

Hence the $T_4$ state is selected $\mathcal Q_+$. The selected-chord estimate
\zcref{lem:new-selected-chords} and the supplied floor $B_4\ge h_0$
give \zcref{eq:ce1-first-chord}.
\end{proof}

\begin{proposition}[CE1 tail-input three-transverse obstruction]
\label{prop:new-ce1-direct-certificate}
Under the domain, path, and tail hypotheses of
\zcref{lem:ce1-first-return-step}, the additional own-radial demands
\begin{equation}
 C_4\ge1-A,\qquad C_3\ge1-m,\qquad C_2\ge1-D
 \label{eq:ce1-radial-hypotheses}
\end{equation}
are incompatible.
\end{proposition}


\begin{proof}
Apply \zcref{lem:ce1-first-return-step}. The selected-state condition forces
\begin{equation}
 \boxed{X>\frac12.}
 \label{eq:ce1-X-half}
\end{equation}
To see this, selectedness gives
\[
 h_0\le\epsilon(A)<\frac{2A}{1-2A}.
\]
If $X\le1/2$, the lower CE1 inequality gives
\[
 2Rh_0=\eta+A+D\ge w(R+A).
\]
Consequently
\[
 f_R(A)<0,
 \qquad
 f_R(z)=w(R+z)(1-2z)-4Rz.
\]
The quadratic is strictly concave and $f_R(0)>0$.  If $R\le1/2$, then
\[
 f_R(P)\ge\frac{1-E}{2}
 (6E^3-2E^2-5E+2)>0.
\]
If $R>1/2$, then $A<A_*=w/2-\eta$ and
\[
 f_R(A_*)=\frac{1-E}{2}(E+11R-3ER-5)>0.
\]
Concavity contradicts $f_R(A)<0$ in both cases, proving
\zcref{eq:ce1-X-half}.  Thus
\begin{equation}
 m<\frac w2<\frac14.
 \label{eq:ce1-m-quarter}
\end{equation}

If $B_3\ge1/2$, then $B_3>1-X$ and the preceding immediate contradiction
applies.  Assume $B_3<1/2$.  At $T_3$ the reflected incoming reach is
$B_3>h_0>m$, while its assumed transverse radial reach is at least $1-m$.
If this state is not selected $\mathcal Q_+$, the exact local catalogue gives
\[
 B_2\ge1-B_3>\frac12>1-X,
\]
and the remaining center-free edge again contradicts $T_1$.

Thus both $T_4$ and $T_3$ are selected $\mathcal Q_+$.  A second use of
\zcref{lem:new-selected-chords} gives
\[
 B_2>(2-5m)B_3-(1-5m)m.
\]
Since $2-5m>0$, \zcref{eq:ce1-first-chord} yields
\begin{equation}
 \boxed{
 B_2>L_2,
 \qquad
 L_2=(2-5m)L_1-(1-5m)m.
 }
 \label{eq:ce1-second-chord}
\end{equation}

\begin{figure}[!htbp]
\centering
\resizebox{.92\linewidth}{!}{%
\begin{tikzpicture}[
  x=1cm,y=1cm,>=Latex,
  state/.style={draw=black!38,rounded corners=1.5pt,fill=black!2,
    minimum height=.78cm,text width=2.35cm,align=center,font=\scriptsize},
  hard/.style={draw=RadialTeal,rounded corners=1.5pt,fill=RadialTeal!7,
    minimum height=.78cm,text width=2.35cm,align=center,font=\scriptsize},
  arr/.style={-{Stealth[length=1.8mm]},draw=black!60,line width=.65pt},
  harr/.style={-{Stealth[length=1.8mm]},draw=RadialTeal,line width=.85pt}
]
  \node[panellabel,anchor=west] at (-.1,2.55)
    {CE1 reverse path: from a deficient terminal to the final contradiction};

  \node[state] (a0) at (0,1.35) {$A_0\le1-h_0$};
  \node[state] (b5) at (3.05,1.35) {$B_5\ge h_0$};
  \node[state] (b4) at (6.10,1.35) {$B_4\ge h_0$};
  \node[hard]  (b3) at (9.15,1.35) {$B_3>L_1$};

  \node[hard]  (b2) at (0,-.55) {$B_2>L_2>\epsilon(D)$};
  \node[hard]  (a2) at (3.05,-.55) {$A_2\le\epsilon(D)$};
  \node[state] (b1) at (6.10,-.55) {$B_1>1-X$};
  \node[state,draw=DemandRed,fill=DemandRed!7]
    (contra) at (9.15,-.55)
    {$A_1\ge X$, $A_1+B_1\le1$\\contradiction};

  \draw[arr] (a0)--node[above=3pt,figlabel,fill=white,inner sep=.6pt] {companion edge} (b5);
  \draw[arr] (b5)--node[above=3pt,figlabel,fill=white,inner sep=.6pt] {Vd0 monotonicity} (b4);
  \draw[harr] (b4)--node[above=3pt,figlabel,fill=white,inner sep=.6pt] {selected $Q_+$ at $T_4$} (b3);
  \draw[harr] (b3.east)--++(.55,0)--++(0,-1.10)--
    node[right=2pt,figlabel,align=left,fill=white,inner sep=.6pt] {selected $Q_+$ at $T_3$\\and the scalar estimate}
    ++(-10.25,0)--(b2.west);
  \draw[harr] (b2)--node[above=3pt,figlabel,fill=white,inner sep=.6pt] {threshold at $T_2$} (a2);
  \draw[arr] (a2)--node[above=3pt,figlabel,fill=white,inner sep=.6pt] {coverage of $e_{1,2}$} (b1);
  \draw[arr,draw=DemandRed] (b1)--(contra);

  \node[figlabel,atlasbox,align=left,anchor=west] at (1.15,-1.68)
    {grey arrows: edge coverage or nonsupercriticality;\qquad
     teal arrows: the genuinely local high-radial steps};
  \node[figlabel,align=center] at (5.25,-2.33)
    {$L_1$ and $L_2$ are direct selected-arc lower bounds; no composed reach map is introduced.};
\end{tikzpicture}%
}
\caption{Logical structure of the CE1 direct reverse-path certificate.  The
proof starts from the negation of the desired terminal bound at $T_0$, moves
through the actual boundary reaches, uses two selected-arc estimates and one
high-radial threshold, and returns to $T_1$ with incompatible incoming and
outgoing demands.}
\label{fig:ce1-reverse-path}
\end{figure}


We now prove the scalar estimate
\begin{equation}
 \boxed{
 D<\frac1{10},
 \qquad
 L_2>2D+3D^2>\epsilon(D),
 \qquad
 \epsilon(D)<X.
 }
 \label{eq:ce1-scalar-estimate}
\end{equation}
The $T_4$ selected-state inequality and
\zcref{eq:selected-low-root-small-bound} give
\begin{equation}
 D\le D_h:=A(4R-1)+10RA^2-\eta.
 \label{eq:ce1-Dh}
\end{equation}
If $D\ge1/10$, then $A+wD<P$ gives
\[
 A<\overline A=\frac{w(5R-1)}{10}.
\]
The right side of \zcref{eq:ce1-Dh} is increasing in $A$, while
$\eta>Rw/2$, so
\[
 D<U(R):=\overline A(4R-1)+10R\overline A^2-\frac{Rw}{2}.
\]
Direct expansion gives
\[
 \frac1{10}-U(R)=-\frac R{10}P_4(R),
\]
where
\[
 P_4(R)=25R^4-60R^3+26R^2+22R-14.
\]
For $y=2R-1\in(0,1)$,
\[
 16P_4(R)
 =25y^4-20y^3-106y^2+124y-39
 \le-101y^2+124y-39<0.
\]
Hence $U(R)<1/10$, a contradiction.  This proves $D<1/10$.

Put
\[
 \Psi(D)=L_2-2D-3D^2,
 \qquad
 J(D)=L_2-\frac{23}{10}D.
\]
Since $D<1/10$,
\[
 \Psi(D)-J(D)=3D\left(\frac1{10}-D\right)>0.
\]
For fixed $R,A$, differentiation of \zcref{eq:ce1-second-chord} gives
\[
 J'(D)=\frac{S(R,A)}{10R^2},
\]
where
\[
 S(R,A)=100A^2-(40R+50)A+R(20-23R).
\]
The nonempty interval $P-RA\le D\le D_h$ implies
\[
 A>\frac{4Rw}{25R-4}=:A_*.
\]
On the relevant interval $\partial S/\partial A<-45$, and
\[
 S(R,A)<S(R,A_*)
 =-\frac{Rq_3(R)}{(25R-4)^2},
\]
where
\[
 q_3(R)=8775R^3-14260R^2+7928R-1120.
\]
For $y=2R-1$,
\[
 8q_3(R)=8775y^3-2195y^2+997y+3007>0.
\]
Thus $J$ decreases with $D$, so $J(D)\ge J(D_h)$.

At $D=D_h$, one has $h_0=2A+5A^2$.  If $L_{2,h}$ is the corresponding
value, then
\[
 L_{2,h}-A(1+4R)=\frac{A}{R^2}Q(R,A),
\]
where
\begin{align*}
Q(R,A)={}&100A^3R-40A^2R^2-30A^2R+12AR^2-15AR+5A\\
&-4R^3+5R^2-R.
\end{align*}
This polynomial is concave in $A$ on $0\le A\le Rw/2$.  Its endpoint
values are
\[
 Q(R,0)=Rw(4R-1)>0
\]
and
\[
 Q\left(R,\frac{Rw}{2}\right)
 =\frac{Rw}{2}
 (25R^5-30R^4+20R^3-3R^2-7R+3)>0.
\]
The last positivity follows after putting $y=2R-1$:
\[
 32p(R)=25y^5+65y^4+90y^3+106y^2-35y+5>0.
\]
Hence $L_{2,h}>A(1+4R)$.

Finally $A<P=\eta E$, $A<Rw/2$, and $1/E>1+Rw/2$ give
\[
 \frac{D_h}{A}
 <C_*:=4R-2+5R^2w-\frac{Rw}{2}.
\]
Moreover
\[
 1+4R-\frac{23}{10}C_*
 =\frac{230R^3-253R^2-81R+112}{20}>0.
\]
Thus $J(D_h)>0$, so
\[
 L_2>2D+3D^2.
\]
Substitution in the selected-root quadratic gives
$\epsilon(D)<2D+3D^2$, and the argument
proving $\epsilon(A)<X$ also gives $\epsilon(D)<X$.  This proves
\zcref{eq:ce1-scalar-estimate}.

The assumed transverse radial reach at $T_2$ is at least $1-D$.  Since
$B_2>\epsilon(D)$, \zcref{lem:new-direct-threshold}, applied after
reflection, gives
\[
 A_2\le\epsilon(D).
\]
Coverage of $e_{1,2}$ therefore forces
\[
 B_1\ge1-A_2\ge1-\epsilon(D)>1-X,
\]
contradicting $A_1\ge X$ and $A_1+B_1\le1$. Thus the tail-input
boundary and radial hypotheses are incompatible.
\end{proof}
